\chapter{Prompts utilizados}

\section{Criterio de documentación}

Este anexo queda preparado para recoger los prompts relevantes utilizados durante el desarrollo y la redacción de la memoria. No se incluirán consultas menores de corrección ortográfica, formato o búsqueda rápida si no han influido en una decisión técnica. Cada entrada deberá permitir reconstruir qué se preguntó, qué utilidad tuvo la respuesta y cómo se comprobó antes de incorporarla al trabajo.

El formato toma como referencia el anexo de uso de inteligencia artificial revisado en \path{INFO/UTILES/actividad_1_aprox_10_limpio.pdf_sin_metadatos.pdf}. En lugar de copiar su contenido, se adapta el esquema al TFG para separar con claridad el prompt original, el uso dado a la respuesta y la verificación realizada con documentación, código, datos o pruebas propias.

\section{Plantilla de registro}

La tabla~\ref{tab:plantilla-prompts} define los campos que se rellenarán cuando se incorporen los prompts definitivos.

\begin{table}[H]
\centering
\caption{Plantilla de documentación de prompts utilizados.}
\label{tab:plantilla-prompts}
\footnotesize
\setlength{\tabcolsep}{4pt}
\renewcommand{\arraystretch}{0.96}
\begin{tabular}{L{2.9cm}L{10.4cm}}
\toprule
\textbf{Campo} & \textbf{Contenido esperado}\\
\midrule
Identificador & Código breve del prompt, por ejemplo \texttt{P1}, \texttt{P2} o \texttt{P3}.\\
Objetivo & Motivo de la consulta y apartado afectado.\\
Prompt original & Texto utilizado, corrigiendo solo erratas menores.\\
Respuesta aprovechada & Parte de la respuesta incorporada al razonamiento.\\
Comprobación & Fuentes, datos, código o pruebas usados para verificarla.\\
Uso final & Decisión, texto, prueba o cambio incorporado.\\
\bottomrule
\end{tabular}
\caption*{\footnotesize\textit{\textbf{Fuente.} Elaboración propia a partir del modelo de anexo revisado en \path{INFO/UTILES}.}}
\end{table}

\Needspace{10\baselineskip}
\section{Prompts pendientes de incorporación}

Los prompts definitivos se incorporarán cuando se disponga del listado real usado durante el proyecto. Para cada uno se mantendrá el siguiente formato:

\begin{enumerate}[leftmargin=*]
  \item \textbf{Prompt B.1. Título provisional}

  \textbf{Objetivo.} Pendiente de completar.

  \textbf{Prompt original.}
  \begin{quote}
  \small\ttfamily
  Pendiente de incorporar el texto del prompt.
  \end{quote}

  \textbf{Respuesta aprovechada.} Pendiente de completar.

  \textbf{Comprobación realizada.} Pendiente de completar.

  \textbf{Uso final en el proyecto.} Pendiente de completar.
\end{enumerate}
